% SOURCE_STATEMENT: "Develop and analyze a multitask learning framework for Hamiltonians trained on data generated at different levels of electronic-structure theory. Determine how information should be shared across tasks and when transfer improves accuracy and sample efficiency."
% Preparation and final self-audit: CONTRIBUTING.md prompts C and D; source/audit artifacts retained in the local physics_statements_2026_10_11 batch.
% !TeX program = lualatex
\documentclass[11pt]{article}
\usepackage[margin=25mm]{geometry}
\usepackage{fontspec}
\setmainfont{Latin Modern Roman}
\usepackage{amsmath,amssymb,mathtools}
\usepackage{unicode-math}
\setmathfont{Latin Modern Math}
\usepackage{xurl,hyperref}
\hypersetup{colorlinks=true,urlcolor=blue}
\setcounter{secnumdepth}{0}
\setlength{\parindent}{0pt}
\setlength{\parskip}{5pt}
\setlength{\emergencystretch}{3em}
\newcommand{\sref}[2]{\hyperlink{source:#1:#2}{[S#2]}}
\newcommand{\cataloguescope}[1]{\paragraph{Recorded target.}#1}
\newcommand{\R}{\mathbb R}
\newcommand{\C}{\mathbb C}
\newcommand{\N}{\mathbb N}
\newcommand{\Z}{\mathbb Z}
\newcommand{\Q}{\mathbb Q}
\newcommand{\F}{\mathbb F}
\newcommand{\eps}{\varepsilon}
\begin{document}
\section{Multitask Hamiltonian Learning across Levels of Theory}\label{statement}

\subsection{Definitions and mathematical statement}

\paragraph{Shared notation.}$\mathbb N=\{1,2,\ldots\}$; the standard number fields have their usual meanings. All additional conventions are stated locally.


A learned Hamiltonian is an effective operator whose parameters are fitted to electronic-structure data. Distinct levels of theory supply tasks with different approximations and errors. Multitask learning shares selected parameters or representations while allowing task-dependent outputs. Transfer benefit is assessed against corresponding single-task prediction and sample requirements.

\paragraph{Catalogue wording retained for traceability.}
Develop and analyze a multitask learning framework for Hamiltonians trained on data generated at different levels of electronic-structure theory. Determine how information should be shared across tasks and when transfer improves accuracy and sample efficiency.

\paragraph{Scope and source qualification.}The source gives a learning-framework question, not a universal theorem that sharing always improves every task.
\subsection{Short English statement}
Develop Hamiltonian learning that shares information across electronic-structure approximations and determine when that sharing improves prediction and sample efficiency.
\subsection{Sources}
\begin{itemize}
\item[\hypertarget{source:statement:1}{S1}] ULAM.AI and the UnsolvedMath Contributors, UnsolvedMath, catalogue entry. Catalogue curation under CC BY 4.0; definitions and English exposition added locally.
\url{https://huggingface.co/datasets/ulamai/UnsolvedMath}.
\item[\hypertarget{source:statement:2}{S2}] Mathematical Methods in Quantum Chemistry (2025), Oberwolfach Reports, DOI 10.4171/owr/2025/15, pp. 615--688. Original report and mathematical context. Incorporated report material remains under the source CC BY-SA 4.0 terms unless otherwise noted.
\url{https://doi.org/10.4171/owr/2025/15}.
\end{itemize}
\label{end:statement}
\end{document}
